\documentclass[11pt]{article}
\usepackage[margin=1in]{geometry}
\usepackage{amsmath,amssymb,amsthm}
\usepackage[T1]{fontenc}
\newtheorem{definition}{Definition}
\newtheorem{theorem}{Theorem}
\newtheorem{corollary}{Corollary}
\newcommand{\Qspace}{\mathcal Q}
\newcommand{\LC}{L_C}

\title{Finite Radiative-Completion and Detector Nonselection}
\author{The \AE ther-Flow Research Project}
\date{2026-07-30}

\begin{document}
\maketitle

\begin{abstract}
This draft/control source-stage note executes one bounded P9-T06 radiative
case without importing target wave equations.  On the four-cycle it constructs
an exact family of second-order completions of the registered proposal-only
static constraint.  Two positive-kinetic members have the same static equation
but different damping and modal frequencies.  The same finite record also
admits inequivalent rank-two and rank-three mode projectors and inequivalent
linear readouts.  The current source package therefore selects neither a
physical radiative sector nor physical polarizations, speed, damping,
dispersion, or detector coupling.  This is a scoped current-source
nonselection obstruction, not a global impossibility result.
\end{abstract}

\section{Source-sealed finite record}

Let \(C_4\) be the cycle on four source addresses and let
\[
\LC=
\begin{pmatrix}
2&-1&0&-1\\
-1&2&-1&0\\
0&-1&2&-1\\
-1&0&-1&2
\end{pmatrix}.
\]
The registered P8 finite response equation is
\[
 \alpha\LC h-\beta\LC u=0.
\]
For this task take the already admitted null-source fixture \(u=0\),
\(\alpha=\beta=1\), and background \(h_0=0\).  Component shifts are removed
by the quotient
\[
 \Qspace=\mathbb R^4/\operatorname{span}\{(1,1,1,1)\}.
\]
No physical time, spatial length, causal cone, kinetic law, or detector map is
contained in this static record.

An orthonormal eigenbasis is
\[
\begin{aligned}
 v_0&=\tfrac12(1,1,1,1), &\lambda_0&=0,\\
 v_1&=\tfrac1{\sqrt2}(1,0,-1,0), &\lambda_1&=2,\\
 v_2&=\tfrac1{\sqrt2}(0,1,0,-1), &\lambda_2&=2,\\
 v_3&=\tfrac12(1,-1,1,-1), &\lambda_3&=4 .
\end{aligned}
\]
Thus \(\Qspace=\operatorname{span}\{v_1,v_2,v_3\}\).

\begin{definition}[Proposal-only completion family]
For \(m>0\) and \(\gamma\geq0\), define the bookkeeping evolution
\[
  m\ddot q+2m\gamma\dot q+\LC q=0,
  \qquad q(\tau)\in\Qspace .
  \tag{1}
\]
The dots refer only to the task-local update parameter \(\tau\).  Equation
(1) is source-extension data for a nonselection test; it is not adopted as a
physical source law.
\end{definition}

\begin{theorem}[Static-equation-preserving dynamical counterfamily]
\label{thm:counterfamily}
Every member of (1) has the same stationary equation \(\LC q=0\).  On an
eigenmode with eigenvalue \(\lambda\in\{2,4\}\), the characteristic roots are
\[
  s_\lambda=-\gamma\pm
  \sqrt{\gamma^2-\frac{\lambda}{m}} .
  \tag{2}
\]
In particular, the two positive-kinetic fixtures
\[
  A:(m,\gamma)=(1,0),\qquad
  B:(m,\gamma)=(4,\tfrac14)
\]
share the same static equation but have incompatible update spectra:
\[
\begin{array}{c|cc}
 &\lambda=2&\lambda=4\\ \hline
A&\pm i\sqrt2&\pm2i\\
B&-\tfrac14\pm i\tfrac{\sqrt7}{4}
  &-\tfrac14\pm i\tfrac{\sqrt{15}}4 .
\end{array}
\]
Consequently the static source equation does not select a unique kinetic
coefficient, damping coefficient, or modal frequency law.
\end{theorem}

\begin{proof}
Setting \(\dot q=\ddot q=0\) in (1) gives \(\LC q=0\), independently of
\(m\) and \(\gamma\).  Substitution of \(q=e^{s\tau}v\), with
\(\LC v=\lambda v\), gives
\(m s^2+2m\gamma s+\lambda=0\), whose roots are (2).
The two displayed fixtures follow by direct substitution.
\end{proof}

\begin{corollary}[Exact response disagreement]
For \(q(0)=v_1+v_3\) and \(\dot q(0)=0\), fixture \(A\) gives
\[
 q_A(\tau)=\cos(\sqrt2\,\tau)v_1+\cos(2\tau)v_3 .
\]
Fixture \(B\) gives, for \(\lambda\in\{2,4\}\),
\[
 x_{\lambda,B}(\tau)=e^{-\tau/4}
 \left[
 \cos(\omega_\lambda\tau)
 \frac{1}{4\omega_\lambda}\sin(\omega_\lambda\tau)
 \right],
\quad
\omega_2=\frac{\sqrt7}{4},\quad
\omega_4=\frac{\sqrt{15}}4 .
\]
At \(\tau=1\), the \(v_1\) coefficients are respectively
\(0.155943694765374\) and \(0.795370023998395\), while the \(v_3\)
coefficients are respectively \(-0.416146836547142\) and
\(0.607054849167036\).
\end{corollary}

\section{Mode and detector nonselection}

Two exact projectors commute with \(\LC\):
\[
 \Pi_2=v_1v_1^\mathsf T+v_2v_2^\mathsf T,\qquad
 \Pi_3=I-v_0v_0^\mathsf T.
\]
They have ranks two and three.  The source record contains no predicate that
selects either image as the physical radiative subspace.  Calling the
two-dimensional degenerate eigenspace ``two polarizations'' would therefore
be a target-facing interpretation, not a source-derived result.

Likewise every \(d\in\Qspace\) defines a formal linear readout
\[
 D_d[q]=d^\mathsf Tq.
\]
For the exact fixture above, \(d=v_1\) and \(d=v_3\) return different
histories.  The registered P7 operational-device suite supplies guarded source
tokens but no physical interferometer response, arm geometry, strain map, or
calibration selecting \(d\).

The graph eigenvalues provide a formal dispersion token
\(\omega_\lambda^2=\lambda/m-\gamma^2\).  They do not provide a physical
wavenumber or speed: any conversion requires a source-derived physical length
map and physical clock.  Thus neither a luminal speed nor physical damping or
dispersion follows from (1).

\section{Typed target comparison and exact obstruction}

Only after the source outputs were sealed was the registered comparison
family opened.  The target-side exact-GR benchmark has two tensor
polarizations, luminal vacuum propagation, no target vacuum damping or
dispersion, and a metric detector response.  Those facts are comparison-side
burdens only.  They were not premises of \(\LC\), (1), \(\Pi_2\), \(\Pi_3\),
or \(D_d\), and they selected no parameter, projector, readout, or rerun.

The failed object is a total source-derived map from the current static finite
record to a physical radiative sector, polarization space, causal speed,
damping/dispersion law, and detector coupling.  Its present absence is
recorded as
\[
\texttt{OBST-P9T06-RADIATIVE-DYNAMICS-DETECTOR-NONSELECTION-001}.
\]
The P9-T06 case is therefore \texttt{INCONCLUSIVE}, with secondary label
\texttt{FORMAL\_ANALOGY}; the rank-two degeneracy is not a benchmark pass.

This result has status \texttt{blocked\_adoption\_open\_continuation}.
A materially new source-derived kinetic law, physical clock and length map,
radiative-sector selector, and detector interaction can reopen the burden.
No canonical ontology or source law is adopted or rejected.  No effective
metric, Einstein equation, exact-GR recovery, benchmark promotion, Gate E
verdict, proof, publication, push, global no-go theorem, future-extension
impossibility, or completed derivation follows.

\end{document}
